\section{Implementation}
\label{sec:impl}

\begin{guidance}[Implementation (15 marks)]
Choose \textbf{3 to 5 short code excerpts} that were \textbf{not straightforward}
to write, and explain each: what problem it solves, why you wrote it this way,
and what would go wrong with the obvious alternative. Good candidates:
receiving bytes in an interrupt and handing them to the main loop; the frame
assembler with timeout and oversize handling; non-blocking button debounce;
keeping the GUI responsive; separating unsolicited messages from replies.
Keep each listing short (10--25 lines) and narrow enough for one column; trim
comments and boilerplate, and say that you did. Do not paste whole files.
\end{guidance}

\subsection{\ph{Heading that states what excerpt 1 achieves}}
\ph{Explain Listing~\ref{lst:example}: the problem, your approach, and why it
works. Cite the documents you relied on, e.g. the HAL manual~\cite{um1785} or
the reference manual~\cite{rm0091}. Mention any bug you found and how you fixed it; markers value an
honest account of what was hard.}

\begin{lstlisting}[language=C, caption={\ph{What this excerpt does} (\texttt{\ph{file.c}}).}, label=lst:example]
/* Replace with a short excerpt of your own code. */
void example(void)
{
    /* ... */
}
\end{lstlisting}

\subsection{\ph{Heading for excerpt 2}}
\ph{Same pattern: problem, approach, justification, listing.}

\subsection{\ph{Heading for excerpt 3}}
\ph{Same pattern. Include at least one excerpt from the PC client (e.g. how it
reads replies without blocking the GUI).}
